% Standalone Tessera: 27 Oct 2026 Research Meeting
% Visual style follows the 21 Oct 2026 light Madrid decks.
% Facts are taken from the local arXiv:2604.10180v2 PDF only.
% Compile with pdfLaTeX from this directory. Keep Tessera/ beside this file.
\documentclass[aspectratio=169,10pt]{beamer}
\usepackage[english]{babel}
\usepackage[utf8]{inputenc}
\usepackage[T1]{fontenc}
\usepackage{lmodern}
\usepackage{amsmath,amssymb}
\usepackage{booktabs,tabularx,array}
\usepackage{graphicx}
\graphicspath{{./}{figures/}}
\usepackage{xcolor}
\usepackage{tikz}
\usetikzlibrary{arrows.meta,positioning,calc,shapes.geometric}
\usepackage{hyperref}

\usetheme{Madrid}
\usecolortheme{default}
\definecolor{Ink}{HTML}{243748}
\definecolor{Muted}{HTML}{667B8B}
\definecolor{Accent}{HTML}{347E74}
\definecolor{AccentBlue}{HTML}{507DAD}
\definecolor{SoftMint}{HTML}{E6F4EE}
\definecolor{SoftBlue}{HTML}{E9F2F9}
\definecolor{SoftLilac}{HTML}{F1EEF8}
\definecolor{Paper}{HTML}{FBFCFE}
\definecolor{Divider}{HTML}{D8E4EB}

\setbeamercolor{background canvas}{bg=Paper}
\setbeamercolor{normal text}{fg=Ink}
\setbeamercolor{structure}{fg=Accent}
\setbeamercolor{frametitle}{fg=Ink,bg=SoftBlue}
\setbeamercolor{title}{fg=Ink}
\setbeamercolor{subtitle}{fg=Muted}
\setbeamercolor{block title}{fg=Ink,bg=SoftBlue}
\setbeamercolor{block body}{fg=Ink,bg=white}
\setbeamercolor{item}{fg=Accent}
\setbeamercolor{caption name}{fg=Muted}
\setbeamercolor{footline}{fg=Muted,bg=Paper}
\setbeamerfont{frametitle}{size=\large,series=\bfseries}
\setbeamerfont{footline}{size=\tiny}
\setbeamerfont{block title}{series=\bfseries}
\setbeamertemplate{headline}{}
\setbeamertemplate{navigation symbols}{}
\setbeamertemplate{blocks}[rounded][shadow=false]
\setbeamertemplate{itemize items}[circle]
\setbeamertemplate{itemize subitem}[circle]
\setbeamertemplate{footline}{%
  \leavevmode\hbox{%
    \begin{beamercolorbox}[wd=\paperwidth,ht=.30cm,dp=.16cm,rightskip=.60cm plus 1fil]{footline}%
      \hfill\insertframenumber\,/\,\inserttotalframenumber
    \end{beamercolorbox}}}
\setbeamertemplate{frametitle}{%
  \nointerlineskip
  \begin{beamercolorbox}[wd=\paperwidth,ht=.72cm,dp=.16cm,leftskip=.55cm,rightskip=.40cm]{frametitle}%
    \usebeamerfont{frametitle}\insertframetitle
  \end{beamercolorbox}}

\newcommand{\smalllabel}[1]{{\scriptsize\bfseries\color{Accent}\MakeUppercase{#1}}}
\newcommand{\hint}[1]{{\scriptsize\color{Muted}#1}}
\newcommand{\takeaway}[1]{%
  \begin{beamercolorbox}[rounded=true,sep=0.12cm,wd=\linewidth]{block title}%
    {\small\bfseries #1}
  \end{beamercolorbox}}
\newcommand{\stepcard}[3]{%
  \begin{beamercolorbox}[rounded=true,sep=0.10cm,wd=\linewidth]{block body}%
    {\scriptsize\bfseries\color{Accent}#1}\par\vspace{.4mm}
    {\footnotesize\bfseries #2}\par\vspace{.4mm}
    {\scriptsize\color{Muted}#3}%
  \end{beamercolorbox}}

\newcommand{\figureplaceholder}[3]{%
  \begingroup\centering
  \IfFileExists{#1}{%
    \includegraphics[width=.98\linewidth,height=#2,keepaspectratio]{#1}%
  }{%
    \IfFileExists{figures/#1}{%
      \includegraphics[width=.98\linewidth,height=#2,keepaspectratio]{figures/#1}%
    }{%
      \begin{tikzpicture}
        \node[draw=AccentBlue!55,dashed,line width=.7pt,rounded corners=4pt,
              fill=SoftBlue!60,minimum height=#2,minimum width=.95\linewidth,
              text width=.85\linewidth,align=center,inner sep=2mm] {
           {\color{AccentBlue}\bfseries FIGURE PLACEHOLDER}\\[2.2mm]
           {\color{Ink}\small #3}\\[1.5mm]
           {\color{Muted}\scriptsize\ttfamily\detokenize{#1}}};
      \end{tikzpicture}%
    }%
  }\par\endgroup}

\newcommand{\src}[1]{\vspace{0.05cm}\hint{Source: #1}}


\title[Tessera]{Tessera: Unlocking Heterogeneous GPUs through Kernel Disaggregation}
\subtitle{Offline Dependencies, Pipelined Execution, Online Adaptation}
\author{Yu-Jui Liang}
\institute{National Taiwan University \quad | \quad Academia Sinica}
\date{27 Oct 2026}
\begin{document}
% ----- SLIDE 01 -----
{\setbeamertemplate{footline}{}
\begin{frame}[plain]
  \vspace*{0.72cm}
  \smalllabel{Research meeting}\par
  \vspace{.36cm}
  {\raggedright\color{Ink}\LARGE\bfseries
    Tessera: Unlocking Heterogeneous GPUs\\[0.16cm]
    through Kernel-Granularity Disaggregation\par}
  \vspace{.28cm}
  {\color{Muted}\large Offline Dependencies, Pipelined Execution, Online Adaptation\par}
  \vspace{.72cm}
  {\bfseries Yu-Jui Liang}\par\vspace{.06cm}
  {\small\color{Muted}Advisor: Hsiang-Yun Cheng}\par\vspace{.04cm}
  {\small\color{Muted}National Taiwan University \quad | \quad Academia Sinica \quad | \quad 27 Oct 2026\par}
  \vfill
  {\color{Accent}\rule{\textwidth}{.7pt}}\par\vspace{.08cm}
  {\scriptsize\color{Muted}Paper: Hu et al., arXiv:2604.10180v2 (9 Jul 2026). Local PDF is the sole factual source.}
\end{frame}}

% ----- SLIDE 02 -----
\begin{frame}[t]{Presentation Roadmap}
  \smalllabel{Tessera / Overview}\par
  \vspace{.16cm}
  \begin{enumerate}
    \setlength{\itemsep}{.14cm}
    \item \textbf{Background / Motivation}\par
      {\small\color{Muted}Why phase and block splits still misplace kernels.}
    \item \textbf{Methodology}\par
      {\small\color{Muted}PTX dependencies, a throughput MILP, and a latency policy with a monitor.}
    \item \textbf{Experimental Results}\par
      {\small\color{Muted}Five GPU types, four model families, and clusters up to 16 GPUs.}
    \item \textbf{Discussion / Conclusion}\par
      {\small\color{Muted}Cost, generality, and a comparison with ThunderServe, Hetis, and Helix.}
  \end{enumerate}
  \src{Local Tessera PDF. The comparison page is our synthesis, not a table from the paper.}
\end{frame}

% ----- SLIDE 03 -----
\begin{frame}[t]{Hardware Heterogeneity Is Kernel-Dependent}
  \smalllabel{Tessera / Background}\par
  \vspace{.04cm}
  {\footnotesize\textbf{Table I.} Datasheet rates and relative rental price. Price is normalized to the L40s. Tildes are printed in the paper.}\par
  \vspace{.03cm}
  {\centering\scriptsize
  \begin{tabular}{lrrrrrr}\toprule
    GPU & CUDA FP32 & Tensor BF16 & L2 (MB) & Mem.\ (GB) & BW (GB/s) & Price\\\midrule
    A100 & 19.5 & 312 & 40 & 80 & 1935 & 1.5\\
    H100 & 67 & 989 & 50 & 80 & 3350 & 2.9\\
    B200 & 80 & $\sim$2500 & 126 & 192 & $\sim$8000 & 5.0\\
    L40s & 91.6 & 366.5 & 96 & 48 & 864 & 1\\
    RTX Pro 6000 & 120 & $\sim$500 & 126 & 96 & 1597 & 1.2\\\bottomrule
  \end{tabular}\par}
  \vspace{.05cm}
  \begin{itemize}\footnotesize
    \setlength{\itemsep}{.04cm}
    \item Vs.\ A100, the L40s has \(4.7\times\) the CUDA-core rate and \(1.2\times\) the tensor-core rate, but 45\% of the memory bandwidth.
    \item The paper says L40s and RTX Pro 6000 reach up to \(7.4\times\) higher compute throughput per dollar. It does not spell out the denominator GPU.
  \end{itemize}
  \vspace{.03cm}\takeaway{No GPU in Table I wins compute, bandwidth, capacity, and price at once. These are not serving measurements.}
  \src{Sec.~II-A, Table I. Prices are ``derived mainly'' from Google Cloud standard rates.}
\end{frame}

% ----- SLIDE 04 -----
\begin{frame}[t]{Why Phase- or Block-Level Disaggregation Is Too Coarse}
  \smalllabel{Tessera / Motivation}\par
  \vspace{.02cm}
  \begin{columns}[c,onlytextwidth]
    \begin{column}{.46\textwidth}
      {\footnotesize\begin{itemize}\setlength{\itemsep}{.06cm}
        \item Normal execution runs K1--K6 on one GPU. Control flow is the black order.
        \item Disaggregation may put K1, K3, K4 on GPU A and K2, K5, K6 on GPU B.
        \item Red arrows are data dependencies that used to be local memory reads.
        \item This drawing is a schematic. It is not a profiled trace.
      \end{itemize}}
      \vspace{.06cm}\takeaway{Splitting kernels is useful only if those red edges stay correct.}
    \end{column}
    \begin{column}{.52\textwidth}
      \figureplaceholder{Tessera/Tessera_Figure_01_Kernel_Disaggregation.png}{5.05cm}{Figure 1}
    \end{column}
  \end{columns}
  \src{Sec.~I, Fig.~1. PDF p.~2.}
\end{frame}

% ----- SLIDE 05 -----
\begin{frame}[t]{Kernel as the Scheduling Granularity}
  \smalllabel{Tessera / Motivation}\par
  \vspace{.12cm}
  \begin{columns}[T,onlytextwidth]
    \begin{column}{.32\textwidth}
      \stepcard{TOO FINE}{Thread blocks}{Inside one kernel, blocks are largely homogeneous. Further splitting does not track a different hardware preference.}
    \end{column}
    \begin{column}{.32\textwidth}
      \stepcard{TOO COARSE}{Phase or block}{Prefill vs.\ decode, or attention vs.\ FFN, still mixes kernels that prefer different GPUs.}
    \end{column}
    \begin{column}{.32\textwidth}
      \stepcard{THE UNIT}{One GPU kernel}{The kernel is the user-level launch Tessera schedules. It is not a new model-layer abstraction.}
    \end{column}
  \end{columns}
  \vspace{.18cm}
  \begin{itemize}\footnotesize
    \setlength{\itemsep}{.06cm}
    \item Phase and block methods are also tied to a model family. Mamba has no attention block; diffusion has no prefill/decode split.
    \item Tessera's claim is that the kernel exists in all four families they evaluate.
  \end{itemize}
  \vspace{.08cm}\takeaway{The granularity argument is architectural and empirical. The counts are on the next two slides.}
  \src{Sec.~I and Sec.~II-C.}
\end{frame}

% ----- SLIDE 06 -----
\begin{frame}[t]{Different Kernels Prefer Different GPUs}
  \smalllabel{Tessera / Motivation}\par
  \vspace{.02cm}
  \begin{columns}[c,onlytextwidth]
    \begin{column}{.48\textwidth}
      {\footnotesize\begin{itemize}\setlength{\itemsep}{.06cm}
        \item Profiled pair: A100 and L40s. Models: Llama3 8B, GPT-oss 20B, Mamba 7B, Qwen2.5-VL 7B, SD3.5.
        \item (a) By kernel count, 67\% of kernels are faster on the L40s, averaged across models and tasks.
        \item (b) Weighted by time on the A100, those kernels are 36\% of end-to-end time, and up to 53\% for diffusion.
        \item Count and time answer different questions. A numerous short kernel is not 53\% of the job.
      \end{itemize}}
      \vspace{.04cm}\takeaway{There is a real L40s-favoring slice. It is not ``most of the runtime'' on every model.}
    \end{column}
    \begin{column}{.50\textwidth}
      \figureplaceholder{Tessera/Tessera_Figure_02_Kernel_Heterogeneity.png}{4.85cm}{Figure 2}
    \end{column}
  \end{columns}
  \src{Sec.~II-B, Fig.~2. The motivation text says Mamba 7B; the eval model is Mamba-Codestral 7B.}
\end{frame}

% ----- SLIDE 07 -----
\begin{frame}[t]{Phase Labels Hide Kernel-Level Variation}
  \smalllabel{Tessera / Motivation}\par
  \vspace{.02cm}
  \begin{columns}[c,onlytextwidth]
    \begin{column}{.47\textwidth}
      {\footnotesize\begin{itemize}\setlength{\itemsep}{.05cm}
        \item GPT-oss 20B. Points above the line are faster on the A100; points below are faster on the L40s.
        \item Inside one phase: 45\% of prefill kernels and 57\% of decode kernels are faster on the L40s.
        \item cublasGemv is about 1 FLOP/byte. A100's HBM is \(\sim 2\times\) L40s, and the kernel is \(1.9\times\) faster on A100.
        \item FlashAttention is compute-heavy. L40s is up to \(2.1\times\) faster. Inside blocks: 71\% of attention kernels and 33\% of FFN kernels prefer the L40s.
      \end{itemize}}
    \end{column}
    \begin{column}{.51\textwidth}
      \figureplaceholder{Tessera/Tessera_Figure_03_Kernel_Duration_Ratios.png}{4.95cm}{Figure 3}
    \end{column}
  \end{columns}
  \vspace{.02cm}\takeaway{A phase label is not a hardware preference. Both phases contain both kinds of kernel.}
  \src{Sec.~II-C, Fig.~3. PDF p.~3. Ratios are L40s duration over A100 duration.}
\end{frame}

% ----- SLIDE 08 -----
\begin{frame}[t]{Three Challenges for Kernel Disaggregation}
  \smalllabel{Tessera / Problem}\par
  \vspace{.12cm}
  \begin{columns}[T,onlytextwidth]
    \begin{column}{.32\textwidth}
      \stepcard{1 / CORRECTNESS}{RAW across GPUs}{A kernel on one GPU may read a buffer written on another. Opaque and JIT kernels do not expose that buffer.}
    \end{column}
    \begin{column}{.32\textwidth}
      \stepcard{2 / BUBBLES}{Wait on the dependency}{The downstream GPU idles unless another request's ready kernels cover the transfer.}
    \end{column}
    \begin{column}{.32\textwidth}
      \stepcard{3 / POLICY}{Two objectives}{Offline jobs want throughput. Online serving wants per-request latency, and the load changes.}
    \end{column}
  \end{columns}
  \vspace{.18cm}
  \begin{itemize}\footnotesize
    \setlength{\itemsep}{.06cm}
    \item A bad placement can spend more on communication and imbalance than it saves in kernel time.
    \item Tessera's answer is one analyzer, two MILP objectives, a pipeline, and a windowed monitor.
  \end{itemize}
  \vspace{.08cm}\takeaway{The rest of the talk is those three answers, in that order.}
  \src{Sec.~I. The three challenges are the paper's own framing.}
\end{frame}

% ----- SLIDE 09 -----
\begin{frame}[t]{Tessera System Overview}
  \smalllabel{Tessera / System}\par
  \vspace{.02cm}
  \begin{columns}[c,onlytextwidth]
    \begin{column}{.48\textwidth}
      {\footnotesize\begin{itemize}\setlength{\itemsep}{.06cm}
        \item Analyzer: memory ranges, RAW edges, and per-GPU kernel latency.
        \item Planner: a kernel-to-GPU policy for each dependency graph.
        \item One worker per GPU runs only its kernels and the cut-edge transfers.
        \item The monitor feeds runtime latency back to policy selection.
      \end{itemize}}
      \vspace{.04cm}\takeaway{NVLink and RDMA in the figure are the design picture. Measured links are stated later.}
    \end{column}
    \begin{column}{.50\textwidth}
      \figureplaceholder{Tessera/Tessera_Figure_04_System_Overview.png}{4.95cm}{Figure 4}
    \end{column}
  \end{columns}
  \src{Sec.~III, Fig.~4. PDF p.~4.}
\end{frame}

% ----- SLIDE 10 -----
\begin{frame}[t]{Offline Planning and Online Adaptation}
  \smalllabel{Tessera / System}\par
  \vspace{.12cm}
  \begin{columns}[T,onlytextwidth]
    \begin{column}{.49\textwidth}
      \stepcard{OFFLINE}{Analyzer and planner}{Run once, usually during warm-up. Build one dependency graph per CUDA graph and solve its MILP. No per-token replanning.}
      \vspace{.14cm}
      \stepcard{ONLINE}{Workers}{Select the precomputed policy when a CUDA graph is launched. Overlap send/recv with compute on separate streams.}
    \end{column}
    \begin{column}{.49\textwidth}
      \stepcard{ONLINE}{Monitor}{Every window \(W\), compare queueing pressure with a threshold and switch between the two policies.}
      \vspace{.14cm}
      \stepcard{NOT IN THE MILP}{GPU memory}{The formulation does not constrain capacity. Weights are replicated, as in prefill/decode disaggregation.}
    \end{column}
  \end{columns}
  \vspace{.16cm}\takeaway{The expensive graph and MILP stay off the request path. The monitor only chooses which stored policy to run.}
  \src{Sec.~III and Sec.~IV. Memory replication is stated in Sec.~IV.}
\end{frame}

% ----- SLIDE 11 -----
\begin{frame}[t]{Analyzer: Standard vs.\ Opaque Kernels}
  \smalllabel{Tessera / Analyzer}\par
  \vspace{.10cm}
  \begin{columns}[T,onlytextwidth]
    \begin{column}{.49\textwidth}
      \stepcard{STANDARD}{Library kernels}{cublasSgemm reads A and B and writes C. Sizes come from the API arguments. No instrumentation is required to know the buffers.}
      \vspace{.12cm}
      \stepcard{OPAQUE}{User or JIT kernels}{Semantics are not available when the framework launches them. Tessera instruments PTX instead of guessing API arguments.}
    \end{column}
    \begin{column}{.49\textwidth}
      \stepcard{WHY NOT API GUESSING}{Caching allocator}{PhoenixOS can see base addresses, but PyTorch's allocator does not call CUDA alloc on every tensor. The guessed size is often too large.}
      \vspace{.12cm}
      \stepcard{COST OF A LARGE GUESS}{Extra traffic}{An overestimated buffer is still transferred. Tessera keeps only the addresses the kernel actually touches.}
    \end{column}
  \end{columns}
  \vspace{.14cm}\takeaway{Standard kernels are read from their API. Opaque kernels are read from instrumented PTX.}
  \src{Sec.~III-A. PhoenixOS is cited as prior buffer discovery, not as a Tessera component.}
\end{frame}

% ----- SLIDE 12 -----
\begin{frame}[t]{PTX Instrumentation and Address Ranges}
  \smalllabel{Tessera / Analyzer}\par
  \vspace{.08cm}
  \begin{enumerate}\footnotesize
    \setlength{\itemsep}{.07cm}
    \item Add an instrumentation buffer to the kernel parameters and load it with \texttt{ld.param}.
    \item Around each \texttt{ld.global} / \texttt{st.global}, emit \texttt{atom.global.min} and \texttt{atom.global.max} so threads agree on one range.
    \item After the launch, align each instruction's \([\min,\max]\) interval with buffer base addresses. The opcode says read or write.
    \item Recompile the patched PTX with \texttt{ptxas}. The probe engine is built on NEUTRINO.
  \end{enumerate}
  \vspace{.06cm}
  {\footnotesize If a kernel touches a global that is not a launch argument, the base address is unknown. Tessera then refuses to disaggregate that kernel. The paper cites a study that this pattern is absent in PyTorch and vLLM, and present in one Rodinia kernel.}\par
  \vspace{.05cm}\takeaway{The analyzer records touched address ranges. It does not reconstruct a kernel's arithmetic.}
  \src{Sec.~III-A and Sec.~IV. Listing 1 is the paper's instrumentation sketch.}
\end{frame}

% ----- SLIDE 13 -----
\begin{frame}[t]{Extracting RAW Dependencies}
  \smalllabel{Tessera / Analyzer}\par
  \vspace{.10cm}
  \begin{itemize}\footnotesize
    \setlength{\itemsep}{.08cm}
    \item Execution order is preserved, so only Read-After-Write edges can change the result. WAR and WAW are not the edges Tessera schedules.
    \item A registry stores the last writer of each buffer. A later reader adds an edge from that writer, annotated with the measured byte size.
    \item Asynchronous copies are nodes too. Their source, destination, and size come from the API.
    \item One CUDA graph gets one dependency graph. Dynamic shapes become several graphs, not one graph with a runtime branch.
    \item Frameworks without CUDA Graph can record a trace and replay it under stream capture. That tool is preprocessing, not the online scheduler.
  \end{itemize}
  \vspace{.06cm}\takeaway{The graph is a static RAW graph per captured execution, plus the byte size of each edge.}
  \src{Sec.~III-A. ``DDG'' is the paper's term for this graph.}
\end{frame}

% ----- SLIDE 14 -----
\begin{frame}[t]{Cross-Iteration Dependencies and KV State}
  \smalllabel{Tessera / Analyzer}\par
  \vspace{.10cm}
  \begin{itemize}\footnotesize
    \setlength{\itemsep}{.08cm}
    \item The default graph is intra-iteration. KV cache is written in one iteration and read in a later one.
    \item The analyzer flags those buffers by profiling more than one iteration.
    \item At runtime each GPU keeps a replica. Only the delta is copied, asynchronously.
    \item Two NCCL communicators separate this background copy from the latency-critical cut-edge transfers. Dependency traffic gets the higher class.
    \item The paper says delta replication is under 0.1\% of latency, citing prior prefill/decode work rather than a new Tessera measurement.
  \end{itemize}
  \vspace{.06cm}\takeaway{KV state is replicated and updated in the background. It is not a cut edge on the per-iteration critical path.}
  \src{Sec.~III-A and Sec.~IV. The 0.1\% figure is attributed to prior work [9], [10].}
\end{frame}

% ----- SLIDE 15 -----
\begin{frame}[t]{Planner Inputs, Decisions, and Objectives}
  \smalllabel{Tessera / Planner}\par
  \vspace{.06cm}
  {\centering\footnotesize
  \begin{tabular}{ll}\toprule
    Input & Meaning\\\midrule
    \(G\), \(K\), \(E\) & GPUs, kernels, and RAW edges\\
    \(t_{k,g}\) & Profiled latency of kernel \(k\) on GPU \(g\)\\
    \(d_{ij}\) & Bytes on the edge from kernel \(i\) to \(j\)\\
    \(bw_{u,g}\), \(\ell_{u,g}\) & Bandwidth and base latency between GPUs \(u\) and \(g\)\\\bottomrule
  \end{tabular}\par}
  \vspace{.08cm}
  \begin{itemize}\footnotesize
    \setlength{\itemsep}{.05cm}
    \item Decision: binary \(x_{k,g}\). Each kernel is placed on exactly one GPU.
    \item Throughput objective, for offline batches: minimize the slowest GPU's stage time, assuming communication can overlap computation.
    \item Latency objective, for a lightly loaded online request: minimize kernel time plus transfer time along the cut, with no overlap credit.
  \end{itemize}
  \vspace{.04cm}\takeaway{Same graph, two objectives. The monitor chooses which solved policy to run.}
  \src{Sec.~III-B, Table II. Symbols are the paper's; the overlap remark is the surrounding prose.}
\end{frame}

% ----- SLIDE 16 -----
\begin{frame}[t]{Offline Throughput-Oriented MILP}
  \smalllabel{Tessera / Planner}\par
  \vspace{.08cm}
  \[
    \sum_{g\in G} x_{k,g}=1,
    \qquad
    T_g=\sum_{k\in K} t_{k,g}\,x_{k,g},
    \qquad
    W_g=\max(T_g,M_g),
    \qquad
    \min \max_{g\in G} W_g.
  \]
  \begin{itemize}\footnotesize
    \setlength{\itemsep}{.06cm}
    \item \(T_g\) is compute time on GPU \(g\). \(M_g\), on the next slide, is communication time into \(g\).
    \item \(W_g\) is the slower of the two. Steady-state throughput is then one request per \(\max_g W_g\).
    \item This is the offline objective. It is the wrong objective when too few requests are in flight to hide transfers.
    \item Gurobi solves it offline. A reported case, GPT-oss 20B on A100+L40s, has about 500 kernels and finishes one graph in about 20 ms.
  \end{itemize}
  \vspace{.04cm}\takeaway{The MILP balances GPUs by minimizing the slowest stage. It does not maximize a sum of kernel speeds.}
  \src{Sec.~III-B. The 20 ms figure is for \(|G|=2\), \(|K|\approx 500\); all patterns take about 5 s.}
\end{frame}

% ----- SLIDE 17 -----
\begin{frame}[t]{Communication and Load-Balancing Constraints}
  \smalllabel{Tessera / Planner}\par
  \vspace{.06cm}
  \[
    c^{u,g}_{ij}=\ell_{u,g}+\frac{d_{ij}}{bw_{u,g}},
    \qquad
    M_g=\sum_{(i,j)\in E}\sum_{u\in G\setminus\{g\}} c^{u,g}_{ij}\,y^{u,g}_{ij}.
  \]
  \begin{itemize}\footnotesize
    \setlength{\itemsep}{.06cm}
    \item \(y^{u,g}_{ij}=x_{i,u}x_{j,g}\) is 1 when the producer is on \(u\) and the consumer is on \(g\). The product is linearized with the standard Boolean replacement.
    \item Load balance is not a separate variance constraint. Minimizing \(\max_g W_g\) already penalizes the busiest GPU.
    \item Per-GPU memory capacity is absent. At initialization every GPU receives a full copy of the weights.
    \item Offline, once the policy is fixed, unused weight buffers can be reclaimed. The paper does not report that reclamation in the evaluation tables.
  \end{itemize}
  \vspace{.04cm}\takeaway{Communication is bytes over profiled bandwidth, plus a base latency. Memory is a runtime caveat, not a constraint.}
  \src{Sec.~III-B and Sec.~IV. Linearization cites Nemhauser and Wolsey.}
\end{frame}

% ----- SLIDE 18 -----
\begin{frame}[t]{Online Latency-Oriented Scheduling}
  \smalllabel{Tessera / Planner}\par
  \vspace{.08cm}
  {\footnotesize Under low load the pipeline cannot hide transfers. The objective adds kernel time and cut-edge transfer time:}\par
  \[
    \min \sum_{k,g} t_{k,g}x_{k,g}
    +\sum_{(i,j)}\sum_{u}\sum_{g\neq u} c^{u,g}_{ij}\,y^{u,g}_{ij}.
  \]
  \begin{itemize}\footnotesize
    \setlength{\itemsep}{.05cm}
    \item There is no \(\max(T,M)\) here. Overlap is not assumed.
    \item The same Gurobi solve is offline. Workers only look up the stored policy.
    \item If communication is expensive, this objective keeps dependent kernels on the same GPU. Fig.~11 shows that behavior at 25 Gb/s.
  \end{itemize}
  \vspace{.04cm}\takeaway{Latency mode shortens one request's critical path. It is not the throughput MILP with a different label.}
  \src{Sec.~III-B. The sum is the paper's low-load decomposition, compacted for the slide.}
\end{frame}

% ----- SLIDE 19 -----
\begin{frame}[t]{Pipelining Across Requests}
  \smalllabel{Tessera / Execution}\par
  \vspace{.10cm}
  \begin{itemize}\footnotesize
    \setlength{\itemsep}{.08cm}
    \item Without a pipeline, a GPU that waits on a remote buffer sits idle. The paper's later drill-down finds about half the optimal throughput in that mode.
    \item Each in-flight request gets its own CUDA stream. A blocked stream does not stop another request's ready kernels.
    \item Equal-priority streams tend to reach their transfers together, so the GPU still goes idle.
    \item Later requests get a lower CUDA stream priority. The hardware scheduler then prefers earlier requests and staggers the transfers.
    \item ``Optimal'' in that drill-down means zero communication overhead, not a better kernel algorithm.
  \end{itemize}
  \vspace{.06cm}\takeaway{The pipeline hides transfers only when several requests are in flight and they are not synchronized.}
  \src{Sec.~III-C. The half-throughput and priority results are measured in Sec.~V-D, Fig.~9.}
\end{frame}

% ----- SLIDE 20 -----
\begin{frame}[t]{GPU Worker Execution}
  \smalllabel{Tessera / Execution}\par
  \vspace{.02cm}
  \begin{columns}[c,onlytextwidth]
    \begin{column}{.48\textwidth}
      {\footnotesize\begin{itemize}\setlength{\itemsep}{.05cm}
        \item Fig.~5 is GPU B from the six-kernel example. Compute and communication use different streams.
        \item \(R_{ij}\) receives kernel \(i\)'s output for kernel \(j\). \(S_{ij}\) sends it. CUDA events order the two streams.
        \item Receives are posted before the consumer kernel. Sends are posted after the producer if the edge is cut.
        \item Sends and receives are GPU-side kernels via NCCL's GPUDirect Async (IBGDA), so the CPU is not on the critical path.
      \end{itemize}}
      \vspace{.04cm}\takeaway{A worker runs its own kernels and the cut-edge messages. It does not run the whole model.}
    \end{column}
    \begin{column}{.50\textwidth}
      \figureplaceholder{Tessera/Tessera_Figure_05_GPU_Worker_Execution.png}{4.15cm}{Figure 5}
    \end{column}
  \end{columns}
  \src{Sec.~III-C, Fig.~5. PDF p.~6.}
\end{frame}

% ----- SLIDE 21 -----
\begin{frame}[t]{Overlapping Communication and Computation}
  \smalllabel{Tessera / Execution}\par
  \vspace{.10cm}
  \begin{itemize}\footnotesize
    \setlength{\itemsep}{.08cm}
    \item Tessera intercepts CUDA Graph creation, partitions nodes onto GPUs, and embeds send/recv kernels in each subgraph.
    \item Events that order the compute and communication streams become ordinary edges inside that subgraph.
    \item Each worker keeps a table from the original graph handle to the policy and the subgraphs. Launch picks the row; it does not rebuild the graph.
    \item Collectives of an existing tensor-parallel group stay on the original homogeneous GPUs. Tessera only reschedules compute kernels inside a heterogeneous pair.
    \item That composition is how the 8 B200 + 8 H100 experiment is run. It is not a claim that every collective was rewritten.
  \end{itemize}
  \vspace{.05cm}\takeaway{Graph replay is the fast path. The MILP has already finished before the first replay.}
  \src{Sec.~III-C and Sec.~IV. TP composition is also used in Sec.~V-C.}
\end{frame}

% ----- SLIDE 22 -----
\begin{frame}[t]{Online Monitor Signals}
  \smalllabel{Tessera / Runtime}\par
  \vspace{.10cm}
  \begin{itemize}\footnotesize
    \setlength{\itemsep}{.08cm}
    \item The monitor records per-request latency, communication time, and per-kernel-group latency.
    \item A kernel group is the run of kernels between two consecutive communication operations. Individual kernels are not timed online.
    \item \(\bar L_{\mathrm{req}}\) is average end-to-end request latency. \(\bar L_{\mathrm{exec}}\) sums compute and communication and excludes queueing.
    \item The pressure signal is the ratio \(\bar L_{\mathrm{req}}/\bar L_{\mathrm{exec}}\).
    \item A low ratio means little queueing, so the latency policy is preferred. A high ratio means queueing dominates, so the throughput policy is preferred.
  \end{itemize}
  \vspace{.06cm}\takeaway{The switch looks at queueing pressure, not at a hand-named GPU counter.}
  \src{Sec.~III-D. Kernel group is the paper's definition.}
\end{frame}

% ----- SLIDE 23 -----
\begin{frame}[t]{Policy Switching and Sensitivity Parameters}
  \smalllabel{Tessera / Runtime}\par
  \vspace{.10cm}
  \begin{columns}[T,onlytextwidth]
    \begin{column}{.49\textwidth}
      \stepcard{WINDOW \(W\)}{How often to decide}{The ratio is computed at window boundaries. The chosen defaults are \(W=300\) ms and \(\beta=1.5\).}
      \vspace{.12cm}
      \stepcard{THRESHOLD \(\beta\)}{When queueing wins}{Crossing \(\beta\) selects the throughput policy. Staying below it selects the latency policy.}
    \end{column}
    \begin{column}{.49\textwidth}
      \stepcard{SWITCH COST}{About 30 ms}{Workers synchronize on an iteration boundary. The 30 ms is the paper's reported stall, not a bound.}
      \vspace{.12cm}
      \stepcard{WHERE IT WAS TUNED}{One trace}{GPT-oss 20B, 360 s of the Azure Conversation trace. Fig.~10 is the sensitivity plot.}
    \end{column}
  \end{columns}
  \vspace{.14cm}\takeaway{\(W\) and \(\beta\) are the only two knobs the paper names. Aggressive and lazy settings both increase latency.}
  \src{Sec.~III-D and Sec.~V-D. Numeric penalties are on the sensitivity slide.}
\end{frame}

% ----- SLIDE 24 -----
\begin{frame}[t]{Implementation Stack}
  \smalllabel{Tessera / Implementation}\par
  \vspace{.12cm}
  \begin{columns}[T,onlytextwidth]
    \begin{column}{.49\textwidth}
      \stepcard{GPU WORKER}{About 3k lines}{C++ and CUDA. Integrated with vLLM and PyTorch. Evaluation uses vLLM v0.18.0.}
      \vspace{.14cm}
      \stepcard{ANALYZER}{About 2k lines}{Python. PTX injection, RAW construction, and the preprocessing capture tool.}
    \end{column}
    \begin{column}{.49\textwidth}
      \stepcard{SOLVER}{Gurobi, offline}{One graph at \(|G|=2\), \(|K|\approx 500\): about 20 ms. Enumerating patterns: about 5 s. Not on the request path.}
      \vspace{.14cm}
      \stepcard{FALLBACK}{No disaggregation}{Indirect global accesses that are not launch arguments stay on one GPU. The paper treats this as rare for PyTorch and vLLM.}
    \end{column}
  \end{columns}
  \vspace{.14cm}\takeaway{The line counts are the paper's. They are not a complexity comparison with Helix.}
  \src{Sec.~IV and the solver paragraph in Sec.~III-B.}
\end{frame}

% ----- SLIDE 25 -----
\begin{frame}[t]{Communication Backends and Interconnects}
  \smalllabel{Tessera / Implementation}\par
  \vspace{.10cm}
  \begin{itemize}\footnotesize
    \setlength{\itemsep}{.07cm}
    \item The communication backend in the evaluation is NCCL. Cut-edge sends use GPU-initiated IBGDA.
    \item Local pairs: A100+L40s, H100+RTX Pro 6000, and B200+H100. Each GPU has a 200 Gb/s RDMA NIC. The paper does not say these pairs use NVLink.
    \item Cluster, inside a homogeneous high-end group: NVLink. Across nodes: 200 Gb/s for 2 A100 + 1 L40s, and 400 Gb/s for 8 B200 + 8 H100.
    \item Fig.~4 draws NVLink between high-end GPUs and RDMA toward general-purpose GPUs. That is the design diagram, not a per-experiment wiring list.
    \item The paper explicitly does not evaluate AMD GPUs, NPUs, PCIe-only hosts, or Ethernet.
  \end{itemize}
  \vspace{.06cm}\takeaway{Say NVLink only for the homogeneous group inside the cluster experiments.}
  \src{Sec.~IV and Sec.~V-A. ``Respectively'' in the bandwidth sentence maps onto the two cluster setups.}
\end{frame}

% ----- SLIDE 26 -----
\begin{frame}[t]{Evaluation Methodology}
  \smalllabel{Tessera / Evaluation}\par
  \vspace{.04cm}
  {\centering\scriptsize
  \begin{tabular}{lll}\toprule
    Family & Model & Offline data\\\midrule
    LLM & Llama-3 8B, GPT-oss 20B & Splitwise conversation; median 1020 in / 129 out\\
    SSM & Mamba-Codestral 7B & Same conversation trace\\
    MLLM & Qwen2.5-VL 7B & COCO captions, images resized to \(512\times 512\)\\
    Diffusion & Stable Diffusion 3.5 & PartiPrompts, \(1024\times 1024\), 28 denoising steps\\\bottomrule
  \end{tabular}\par}
  \vspace{.05cm}
  \begin{itemize}\footnotesize
    \setlength{\itemsep}{.04cm}
    \item Baselines: homogeneous left and right GPUs; PD follows DistServe; AF follows MegaScale-Infer. PD does not apply to diffusion. AF does not apply to SSM or diffusion.
    \item Engine: vLLM v0.18.0, continuous batching on, DVFS off. Each run is repeated ``multiple times''; variance is below 1\%.
    \item Throughput is output tokens/s for LLM, SSM, and MLLM. Diffusion is images/min on Fig.~6, which the text also calls requests per minute.
    \item Online experiments are GPT-oss 20B only: Poisson arrivals, plus one Azure trace for the monitor.
  \end{itemize}
  \vspace{.03cm}\takeaway{A red X on Fig.~6 means the baseline does not apply. It is not a zero-throughput measurement.}
  \src{Sec.~V-A. Normalized online latency is end-to-end time divided by output tokens.}
\end{frame}

% ----- SLIDE 27 -----
\begin{frame}[t]{Offline Throughput}
  \smalllabel{Tessera / Results}\par
  \vspace{.02cm}
  \figureplaceholder{Tessera/Tessera_Figure_06_Offline_Throughput.png}{3.15cm}{Figure 6}
  {\centering\hint{Fig.~6. Three pairs. Left axis: tokens/s. Right axis, SD only: images/min. Red X: method does not apply.}\par}
  \vspace{.04cm}
  \begin{itemize}\footnotesize
    \setlength{\itemsep}{.04cm}
    \item Against PD disaggregation: \(1.5\times\) on average, up to \(2.3\times\). Against AF disaggregation: \(1.35\times\) on average, up to \(1.7\times\).
    \item The paper does not assign those maxima to one named bar. Do not read a single panel as ``the'' \(2.3\times\).
    \item Gains are attributed to scheduling kernels that a phase or a block would have kept together.
  \end{itemize}
  \vspace{.03cm}\takeaway{The \(2.3\times\) is the largest reported ratio against PD here, not a typical per-request speedup.}
  \src{Sec.~V-B, Fig.~6. PDF p.~8. Five workloads times three pairs; some baseline bars are absent.}
\end{frame}

% ----- SLIDE 28 -----
\begin{frame}[t]{Online Latency under a Token SLO}
  \smalllabel{Tessera / Results}\par
  \vspace{.02cm}
  \begin{columns}[c,onlytextwidth]
    \begin{column}{.48\textwidth}
      {\footnotesize\begin{itemize}\setlength{\itemsep}{.05cm}
        \item GPT-oss 20B, Poisson arrivals, three pairs. Y-axis: end-to-end latency per output token.
        \item The red line is the SLO, 50 ms/token. This is not DistServe's goodput definition.
        \item At low rate, Tessera's normalized latency is \(1.3\times\) and \(1.2\times\) lower than PD and AF.
        \item Under the SLO, Tessera holds up to \(1.3\times\) the request rate of the best baseline.
        \item The abstract's ``\(1.3\times\) more requests under the same SLO'' is this comparison.
      \end{itemize}}
      \vspace{.04cm}\takeaway{Low-rate latency and high-rate SLO capacity are both reported. They are different operating points.}
    \end{column}
    \begin{column}{.50\textwidth}
      \figureplaceholder{Tessera/Tessera_Figure_07_Online_Latency.png}{4.55cm}{Figure 7}
    \end{column}
  \end{columns}
  \src{Sec.~V-B, Fig.~7. PDF p.~8. Online scope is this one LLM.}
\end{frame}

% ----- SLIDE 29 -----
\begin{frame}[t]{Scaling with Cluster Size}
  \smalllabel{Tessera / Results}\par
  \vspace{.02cm}
  \figureplaceholder{Tessera/Tessera_Figure_08_Cluster_Scale_Throughput.png}{2.65cm}{Figure 8}
  {\centering\hint{Fig.~8. Offline tokens/s. (a) 2 A100 + 1 L40s. (b) 8 B200 + 8 H100, kernel split inside each TP rank.}\par}
  \vspace{.03cm}
  {\centering\scriptsize
  \begin{tabular}{lccc}\toprule
    Setup & PD & AF & Tessera\\\midrule
    GPT-oss 20B, 2 A100 + 1 L40s & 2987 & 3144 & 4253\\
    Qwen-3 235B, 8 B200 + 8 H100 & 4367 & 4912 & 6799\\\bottomrule
  \end{tabular}\par}
  \vspace{.03cm}
  {\scriptsize PD on (a) uses the L40s for prefill and each A100 as a decode instance. AF runs attention on the two A100s at TP\,=\,2 and FFN on the L40s. The paper summarizes both setups as \(1.5\times\) vs.\ PD and \(1.4\times\) vs.\ AF.}\par
  \vspace{.02cm}\takeaway{These are offline cluster throughputs. Qwen-3 235B is not in the single-pair Fig.~6 suite.}
  \src{Sec.~V-C, Fig.~8. PDF p.~9. Units are tokens/s.}
\end{frame}

% ----- SLIDE 30 -----
\begin{frame}[t]{Pipeline Effectiveness}
  \smalllabel{Tessera / Drill-down}\par
  \vspace{.02cm}
  \figureplaceholder{Tessera/Tessera_Figure_09_Pipeline_Effectiveness.png}{2.55cm}{Figure 9}
  {\centering\hint{Fig.~9. GPT-oss 20B, offline. The caption does not name the GPU pair.}\par}
  \vspace{.03cm}
  \begin{itemize}\footnotesize
    \setlength{\itemsep}{.03cm}
    \item No pipeline: 1443 tokens/s, about half of the 2734 tokens/s zero-communication reference.
    \item Pipeline: 2123 tokens/s, \(1.47\times\) the no-pipeline bar. Priority: 2641 tokens/s, 96.6\% of that reference.
    \item On the bottleneck GPU, communication idle time falls from nearly half to under 4\%.
    \item 2641 tokens/s matches Fig.~11's 200 Gb/s A100+L40s bar. That match is our cross-figure reading, not a caption statement.
  \end{itemize}
  \vspace{.02cm}\takeaway{Priority matters because it keeps transfers from lining up. The model does not assume the idle time is zero.}
  \src{Sec.~V-D, Fig.~9. PDF p.~9.}
\end{frame}

% ----- SLIDE 31 -----
\begin{frame}[t]{Online Monitor Sensitivity}
  \smalllabel{Tessera / Drill-down}\par
  \vspace{.02cm}
  \begin{columns}[c,onlytextwidth]
    \begin{column}{.48\textwidth}
      {\footnotesize\begin{itemize}\setlength{\itemsep}{.06cm}
        \item GPT-oss 20B, 360 s Azure trace. Left: sweep \(W\) at \(\beta=1.5\). Right: sweep \(\beta\) at \(W=300\) ms.
        \item \(W=30\) ms or \(\beta=1.1\): frequent switches, latency up to \(1.32\times\).
        \item \(W=1.5\) s or \(\beta=3.0\): slow reaction, latency up to \(1.55\times\).
        \item Each switch stalls workers for about 30 ms at an iteration boundary.
        \item The chosen point is \(W=300\) ms, \(\beta=1.5\). It is the best tradeoff in this sweep, not a universal constant.
      \end{itemize}}
    \end{column}
    \begin{column}{.50\textwidth}
      \figureplaceholder{Tessera/Tessera_Figure_10_Online_Monitor_Sensitivity.png}{4.70cm}{Figure 10}
    \end{column}
  \end{columns}
  \vspace{.02cm}\takeaway{Both too many switches and too few switches raise latency on this trace.}
  \src{Sec.~V-D, Fig.~10. PDF p.~10.}
\end{frame}

% ----- SLIDE 32 -----
\begin{frame}[t]{Network Bandwidth and MILP Analysis}
  \smalllabel{Tessera / Drill-down}\par
  \vspace{.02cm}
  \figureplaceholder{Tessera/Tessera_Figure_11_Network_Bandwidth_and_MILP.png}{2.15cm}{Figure 11}
  {\centering\hint{Fig.~11. (a) GPT-oss 20B, A100+L40s. (b) Gurobi time vs.\ kernels per graph.}\par}
  \vspace{.02cm}
  {\footnotesize\textbf{(a)} Throughput at 200 / 100 / 50 / 25 Gb/s: 2641 / 2610 / 2550 / 2490 tokens/s. The 25 Gb/s drop is under 6\%. Normalized latency: 1.00 / 1.03 / 1.08 / 1.15, still below the homogeneous A100 line at 1.18.}\par
  \vspace{.02cm}
  {\footnotesize\textbf{(b)} Text, 2 GPUs: 0.07 s at 200 kernels and 0.43 s at 1000 kernels. A DeepSeek-V3 671B forward graph of about 1500 kernels solves within 1 s. That is a solver stress test, not a serving result. Time grows with the number of GPU types; repeated layers are solved once.}\par
  \vspace{.02cm}\takeaway{At 25 Gb/s the offline pipeline still hides most of the extra delay. The latency policy moves work off the slow cut.}
  \src{Sec.~V-D, Fig.~11. PDF p.~10. Do not read panel (b) as an end-to-end DeepSeek result.}
\end{frame}

% ----- SLIDE 33 -----
\begin{frame}[t]{Cost Efficiency of Heterogeneous Pairing}
  \smalllabel{Tessera / Results}\par
  \vspace{.06cm}
  {\footnotesize\textbf{Table III.} Perf/\$: mean throughput over the five workloads, divided by rental cost, then normalized to the left GPU alone. A pair pays for both GPUs.}\par
  \vspace{.05cm}
  {\centering\footnotesize
  \begin{tabular}{lccc}\toprule
    System & A100+L40s & H100+RTX Pro 6000 & B200+H100\\\midrule
    Homogeneous left & 1.00 & 1.00 & 1.00\\
    Homogeneous right & 1.18 & 2.01 & 0.78\\
    PD disaggregation & 0.87 & 1.00 & 0.70\\
    AF disaggregation & 1.02 & 1.07 & 0.74\\
    Tessera & 1.21 & 1.64 & 1.01\\\bottomrule
  \end{tabular}\par}
  \vspace{.06cm}
  \begin{itemize}\footnotesize
    \setlength{\itemsep}{.04cm}
    \item Vs.\ PD: \(1.5\times\) cost efficiency on average, up to \(1.6\times\). Vs.\ AF: \(1.4\times\) on average, up to \(1.5\times\).
    \item On H100+RTX Pro 6000, Tessera is \(1.64\times\) the Perf/\$ of H100 alone and \(2.3\times\) its throughput, above two H100s, at price \(2.9+1.2\) versus \(2\times 2.9\).
  \end{itemize}
  \vspace{.04cm}\takeaway{The pair can beat a single expensive GPU on Perf/\$. The prices are Table I's relative rentals, not a new quote.}
  \src{Sec.~V-B, Table III. The \(2.3\times\) sentence is the cost-efficiency paragraph, not the PD maximum.}
\end{frame}

% ----- SLIDE 34 -----
\begin{frame}[t]{Generality Across Model Families}
  \smalllabel{Tessera / Results}\par
  \vspace{.06cm}
  {\centering\footnotesize
  \begin{tabularx}{\textwidth}{@{}>{\raggedright\arraybackslash}p{3.3cm}*{2}{>{\raggedright\arraybackslash}X}@{}}\toprule
    What was run & Models & What it does not show\\\midrule
    Offline pairs & Llama-3 8B, GPT-oss 20B, Qwen2.5-VL 7B, Mamba-Codestral 7B, SD 3.5 & PD or AF is inapplicable on some families\\
    Online SLO & GPT-oss 20B only & Latency for SSM, MLLM, or diffusion\\
    Cluster, offline & GPT-oss 20B and Qwen-3 235B & The single-pair Fig.~6 suite\\
    Solver stress test & DeepSeek-V3 671B, one forward graph & End-to-end serving\\\bottomrule
  \end{tabularx}\par}
  \vspace{.08cm}
  \begin{itemize}\footnotesize
    \setlength{\itemsep}{.05cm}
    \item AF needs an attention block, so Mamba and diffusion are out. PD needs a prefill/decode split, so diffusion is out.
    \item The motivation profiles say ``Mamba 7B''; the evaluated SSM is Mamba-Codestral 7B. Do not merge the two names.
  \end{itemize}
  \vspace{.06cm}\takeaway{``Four families'' means the offline pair study. Online evidence is one LLM.}
  \src{Sec.~II-B, Sec.~V-A--V-D. DeepSeek-V3 is the solver paragraph in Sec.~V-D.}
\end{frame}

% ----- SLIDE 35 -----
\begin{frame}[t]{Tessera vs.\ ThunderServe, Hetis, and Helix}
  \smalllabel{Cross-paper comparison}\par
  \vspace{.03cm}
  {\centering\scriptsize
  \renewcommand{\arraystretch}{1.08}
  \begin{tabularx}{\textwidth}{@{}>{\raggedright\arraybackslash}p{1.55cm}*{4}{>{\raggedright\arraybackslash}X}@{}}\toprule
    & ThunderServe & Hetis & Helix & Tessera\\\midrule
    Granularity & Prefill / decode phase & Dense module vs.\ head & Layer range and path & CUDA kernel\\
    Decision & Group, phase, route & Primary set, then heads & Placement, then a path & Kernel-to-GPU map\\
    Time scale & Offline + light reschedule & Offline + online re-dispatch & Once per cluster + per request & Offline MILP + window \(W\)\\
    State moved & KV at phase boundary & QKV and KV shards & Tokens and activations & Cut tensors; KV delta\\
    Adaptation & Phase role, failure & Re-dispatch, migration & KV mask; path reused & Switch on \(\bar L_{\mathrm{req}}/\bar L_{\mathrm{exec}}\)\\
    Environment & Heterogeneous cloud & Local heterogeneous cluster & Heterogeneous, incl.\ geo & NVIDIA pairs; NVLink + RDMA\\
    Models & LLM & LLM & LLaMA 30B / 70B & LLM, SSM, MLLM, diffusion\\
    Main limit & Heuristic phase search & Head groups; local net & One-shot plan; geo simulated & No memory constraint in MILP\\\bottomrule
  \end{tabularx}\par}
  \vspace{.04cm}\takeaway{Kernel splitting is a different cut from phase, head, or layer placement. Do not rank the papers by their largest ratio.}
  \src{Our synthesis. Tessera cells are this PDF. Other cells come from each system's own paper.}
\end{frame}

% ----- SLIDE 36 -----
\begin{frame}[t]{Limitations and Threats to Validity}
  \smalllabel{Tessera / Discussion}\par
  \vspace{.08cm}
  \begin{itemize}\footnotesize
    \setlength{\itemsep}{.06cm}
    \item The MILP ignores per-GPU memory. Full weight replication is conservative, and the paper does not show reclamation in the result tables.
    \item Online latency, the SLO, and the monitor sweep are GPT-oss 20B. Diffusion and Mamba are offline pair results.
    \item The \(<0.1\%\) delta-replication cost is cited from prior PD work. The \(<4\%\) idle time is Tessera's own Fig.~9.
    \item Indirect PTX accesses fall back to one GPU. The paper limits the study to NVIDIA.
    \item Fig.~9's caption does not name the GPU pair. A policy switch stalls about 30 ms.
    \item There is no shared-hardware comparison against Helix, Hetis, or ThunderServe.
  \end{itemize}
  \vspace{.05cm}\takeaway{Generality is real for offline kernels. It is not yet a latency result for every family.}
  \src{Sec.~IV, Sec.~V, and Sec.~VII. The cross-paper gap is ours, not a claim in Tessera.}
\end{frame}

% ----- SLIDE 37 -----
\begin{frame}[t]{Conclusion}
  \smalllabel{Tessera / Conclusion}\par
  \vspace{.12cm}
  \begin{columns}[T,onlytextwidth]
    \begin{column}{.32\textwidth}
      \stepcard{1}{Schedule kernels}{Phase and block labels still mix kernels that prefer different GPUs. The kernel is the unit Tessera places.}
    \end{column}
    \begin{column}{.32\textwidth}
      \stepcard{2}{Make the split correct}{PTX ranges build a RAW graph. The throughput MILP overlaps compute with communication; the latency MILP does not.}
    \end{column}
    \begin{column}{.32\textwidth}
      \stepcard{3}{What was measured}{Up to \(2.3\times\) throughput vs.\ PD and up to \(1.6\times\) Perf/\$. A pair can beat two H100s on the reported H100+RTX setup.}
    \end{column}
  \end{columns}
  \vspace{.16cm}
  \begin{itemize}\footnotesize
    \setlength{\itemsep}{.05cm}
    \item The same system serves models where PD or AF does not apply, in the offline pair study.
    \item The monitor's \(1.3\times\) SLO result is GPT-oss 20B at 50 ms/token, not a goodput number under another paper's definition.
  \end{itemize}
  \vspace{.06cm}\takeaway{Tessera treats GPU diversity as something to schedule per kernel, with an offline plan and a small online switch.}
  \src{Sec.~VII, with ratios from Sec.~V-B rather than from the abstract alone.}
\end{frame}

% ----- SLIDE 38 -----
\begin{frame}
  \vfill
  \begin{center}
    {\fontsize{32}{38}\selectfont\bfseries\color{Ink}Questions \& Discussion}\par
    \vspace{.32cm}
    {\large\color{Accent}Kernel graph \quad / \quad two MILP objectives \quad / \quad what was actually served}\par
    \vspace{.38cm}
    {\small\color{Muted}Tessera (kernel) \quad next to phase, head, and layer-pipeline systems}
  \end{center}
  \vfill
\end{frame}

\end{document}
